% Nature Methods: neutral adaptation of inspected Perspective/Analysis editions.
% No source article text, data, authors, DOI, logo, or original figures.
\PassOptionsToPackage{no-math}{fontspec}
\documentclass[UTF8,fontset=none,scheme=plain]{ctexart}
\usepackage[paperwidth=595.276bp,paperheight=790.866bp,left=39.685bp,right=35.716bp,top=49bp,bottom=43bp,headheight=20bp,headsep=8bp,footskip=24bp]{geometry}
% Shared dependencies and neutral material. Compile with XeLaTeX.
% No publisher logos, paper text, data, figures, or download metadata are included.
\usepackage{amsmath,amssymb,graphicx,xcolor,tikz,fancyhdr,needspace,xurl,hyperref,balance,etoolbox}
% The legacy balance package omits the first split box's final baseline depth.
% Reserve that measured depth rather than hiding a CJK overfull-box warning.
\makeatletter
\patchcmd{\@BAlancecol}{\setbox\@leftcolumn=\vbox to \@tempdima}{\advance\@tempdima by \dp1\relax\setbox\@leftcolumn=\vbox to \@tempdima}{}{\PackageWarning{cns-template}{Balance depth patch unavailable; inspect final columns}}
\makeatother
\usetikzlibrary{arrows.meta,positioning}
\usepackage[protrusion=true,expansion=false]{microtype}
\IfFontExistsTF{Noto Serif CJK SC}{%
  \setCJKmainfont{Noto Serif CJK SC}%
}{%
  \setCJKmainfont[BoldFont=FandolSong-Bold.otf,ItalicFont=FandolKai-Regular.otf]{FandolSong-Regular.otf}%
}
\IfFontExistsTF{Noto Sans CJK SC}{%
  \setCJKsansfont{Noto Sans CJK SC}%
}{%
  \setCJKsansfont[BoldFont=FandolHei-Bold.otf]{FandolHei-Regular.otf}%
}
\setmainfont{TeX Gyre Termes}\setsansfont{TeX Gyre Heros}
\setmonofont{TeX Gyre Cursor}
\definecolor{JournalBlue}{HTML}{0097C9}
\definecolor{JournalRed}{HTML}{B83443}
\definecolor{ScienceRed}{HTML}{CD1021}
\definecolor{CiteBlue}{HTML}{0000EE}
\hypersetup{colorlinks=true,linkcolor=JournalBlue,urlcolor=JournalBlue,
  pdftitle={CNS论文汉化：中性版式示例},pdfauthor={示例模板；非期刊官方模板},
  pdfsubject={Neutral source-inspired Chinese layout demonstration; no real research claims}}
\urlstyle{same}\raggedbottom
\setlength{\parindent}{1em}\setlength{\parskip}{0pt}
\setlength{\emergencystretch}{2em}\tolerance=2200
\XeTeXlinebreaklocale "zh"\XeTeXlinebreakskip=0pt plus .1pt
\pagestyle{fancy}\fancyhf{}
\renewcommand{\headrulewidth}{0pt}\renewcommand{\footrulewidth}{0pt}
\newcommand{\PaperTitle}{中文题名占位：从组织结构到证据解释}
\newcommand{\SampleAuthors}{Author A\textsuperscript{1}, Author B\textsuperscript{2,*}}
\newcommand{\SampleAffiliations}{\textsuperscript{1}作者单位一（示例占位）。\quad
  \textsuperscript{2}作者单位二（示例占位）。\par
  *通信作者信息待填。DOI、收稿和出版日期均应从目标原文核实。}
\newcommand{\SampleAbstract}{本段是中性摘要占位文字，不陈述任何真实研究发现。它演示中文题名、作者信息、摘要和双栏正文之间的层级关系。实际使用时，应完整替换为目标论文的忠实中文译文，并保留研究对象、方法、结果和结论的限定条件。图注、参考文献、作者声明和来源信息均属于独立的覆盖项。模板中的图形仅说明排版结构，不表示生物学机制、实验结果或统计关系。}
\newcommand{\DemoCaption}{\textbf{图1. 中性证据链示意（占位图，无实验数据）。}
  \textbf{A}，输入节点表示需要核对的原文材料。
  \textbf{B}，处理节点表示逐段转换及复核。
  \textbf{C}，输出节点表示正文、图注和书目等独立交付部分。
  箭头仅表示本示例的阅读顺序，不代表因果关系、效应方向或数量变化。
  该图由模板自行绘制，所有节点均为文字占位；无样本量、误差线、统计检验或比例尺。
  实际论文图注应逐项保留原稿报告的样本量、单位、检验方法、限定及图形署名。}
\newcommand{\DemoVector}{%
 \resizebox{\linewidth}{!}{%
 \begin{tikzpicture}[font=\sffamily\small,node distance=18mm,
   box/.style={draw=black!70,rounded corners=2pt,fill=black!3,minimum width=26mm,minimum height=16mm,align=center},
   flow/.style={-{Stealth[length=2.5mm]},line width=.7pt}]
  \node[box] (a) {\textbf{A}\quad 输入\\原文材料};
  \node[box,right=of a] (b) {\textbf{B}\quad 处理\\转换与复核};
  \node[box,right=of b] (c) {\textbf{C}\quad 输出\\中文译稿};
  \draw[flow] (a)--(b);\draw[flow] (b)--(c);
  \node[below=11mm of a,align=center] {标题与摘要\\正文与图表};
  \node[below=11mm of b,align=center] {保留限定\\检查单位};
  \node[below=11mm of c,align=center] {可追溯引文\\明确原文信息};
  \node[below=29mm of b,font=\sffamily\footnotesize,text=black!65] {示意占位；不包含真实实验数据};
 \end{tikzpicture}}%
}
\newcommand{\SampleBibliography}{%
 \section*{参考文献（示例占位）}
 {\fontsize{7.6bp}{10bp}\selectfont\setlength{\parindent}{0pt}\interlinepenalty=10000
 \hypertarget{ref:1}{}1. Author A. (20XX). Example reference title. \textit{Journal placeholder}. 卷、页码与DOI待填。此条仅供演示引文跳转。\par\smallskip
 \hypertarget{ref:2}{}2. Author B. (20YY). Another example reference title. \textit{Journal placeholder}. 出版信息待填；不要把占位条目带入正式译稿。\par}
}
\newcommand{\SampleSourceInfo}{%
 \section*{原文信息}
 以下字段全部为示例占位：作者；英文原标题；期刊；年份；卷期；页码或文章号；DOI或来源链接。实际使用时必须逐项替换并核对。\par
 本文件是\JournalName 风格的非官方中文译稿样板，依据一次已核验文章版式进行结构化改编，不代表该期刊所有年代或文体的统一标准，也不保证像素级复刻。无真实论文内容或期刊标识图像。\par
}
% Deliberately neutral demo prose, not translated from a published paper.
\newcommand{\SampleOpening}{%
 \section*{引言：示例正文与引文}
 这是一段用于检查双栏阅读节奏的中性示例文字。段落中的术语、缩写和标点需要在实际项目中与源文保持一致。中文译稿应呈现原论文的完整论述，而不是把多段证据压缩成一段摘要。这里的引文只链接到本文件末尾明确标为占位的参考文献\SampleCiteOne。\par
 若原句仅表示相关、可能或提示，译文也应保留相同的证据强度。数字、单位、正负号、上下标和范围端点应逐项核对。原文存在内部矛盾时，应分别保留相应位置的陈述，再用简短译注说明；不能根据常识把不同的数字静默统一。\par
 \subsection*{内容层级与跨页衔接}
 此小标题演示二级层级。实际论文可能采用不同的小标题密度，模板不会替代对原文结构的识别。双栏文章需按左栏、右栏和跨页续句的顺序读取。若段落在页末中断，应在完成整句核对后接入下一页，避免把图中文字误当正文，或重复翻译同一段。\par
 图表具有独立的信息范围。正文引用到某幅图，并不意味着其图注可以省略。图注中的每个分面、比较对象、时间点、样本量及统计说明都应对应源文。英文图内标签可以保留，并在图旁或文末提供中文释义；实际采用哪种安排，需要依据图形尺寸和译文长度判断\SampleCiteTwo。\par
 \subsection*{来源与可维护数据}
 本段继续演示正文密度。可编辑文本和图注数据应与最终版面分开维护，以便修订时准确定位变化。若只调整版式，应能证明正文、图注和书目内容没有被意外改写。更换字体或纸张后，须重新检查断行、图注与图片的相对位置以及页脚可读性。\par
 文内引用与书目之间应有明确对应关系。这里的两条演示书目没有真实DOI，也没有虚构的实验信息。使用模板时可增删条目，但必须同时更新编号、链接目标和引用位置。页码只表示中文译稿自己的分页，不用于冒充原出版页码。\par
 \subsection*{结果叙述的演示占位}
 本节的文字仅用于呈现长段落在双栏页面中的流动，不报告真实结果。实际翻译应保留观察对象、实验条件和结局指标之间的关系。例如，原文分别描述体外实验与动物实验时，不能把两类结果合并为同一层面的结论；研究使用的对照、时间顺序和分析范围也必须保留。图表中的比较范围不应扩大为对整个人群或全部疾病阶段的判断。\par
 论述中的转折与限制同样属于正文。若研究先提出一种解释，再讨论另一种可能，译文应沿用这一推进顺序。作者明确写出的不确定性应能在中文中被识别；如果尚未确定某个现象的方向或机制，就不能改成已经建立的因果关系。对于术语近似但研究层级不同的概念，可在首次出现时保留缩写，后文保持统一。\par
 \subsection*{方法信息与可复核性}
 方法部分往往包含密集数字和专有名词。翻译时须区分浓度、总量、频率和处理时长，不能因为中文句式调整而改变参数的归属。试剂名称、软件版本和数据访问标识通常保留原样；上下限、排除规则和显著性标准则需要逐字符核对。这里没有提供任何真实实验参数，避免把模板占位误认为可以直接执行的方法。\par
 表格与补充材料的引用也需要保留原层级。正文对补充图的指向并不意味着模板中已经附有该补充图；若源文件只包含主文，应如实标注所交付的范围。图表编号、附录顺序和参考文献数量都应由目标文章决定，本示例中的结构不构成固定配额，也不要求所有论文采用同样的章节。\par
}
\newcommand{\SampleClosing}{%
 \section*{讨论：版式适配与范围}
 上方矢量图是模板原创的排版占位图，不来自任何论文。图形与完整演示图注被放在同一页，便于读者直接核对分面说明。真实项目应使用经过授权或用户提供的原图素材，并保留图注署名与许可信息。尺寸不足时，可适度缩放或调整栏宽，不以删除图注解决空间问题。\par
 模板仅记录一次已核验文章的可观察设计特征。即使文章出自同一期刊，不同年代、研究论文与综述、校正清样与正式版本也可能采用不同的页眉和引文形式。应以用户提供的实际PDF为准，不能仅凭期刊名称推断配色、字号或引用方式\SampleCiteOne。\par
 \subsection*{演示检查项}
 本样板用于展示可编译性、中文字形、标题层级、跨栏图注、页眉页脚和内部链接。图中没有实际数据，因此不应被当作科学插图引用。替换内容后，还需检查所有页的渲染结果，确认无截字、碰撞、异常空白和遗漏段落。\par
 \SampleBibliography
 \SampleSourceInfo
}


\usepackage{multicol,booktabs,array}
\definecolor{MethodsOrange}{HTML}{F2652E}
\definecolor{MethodsBlue}{HTML}{006AA7}
\definecolor{MethodsBox}{HTML}{EFEEE5}
\hypersetup{colorlinks=true,linkcolor=MethodsBlue,urlcolor=MethodsBlue,
  pdftitle={Nature Methods风格中性中文样板},pdfauthor={非官方中性样板},
  pdfsubject={Original neutral layout demonstration, no source research claims}}
\setlength{\columnsep}{12bp}
\setlength{\parindent}{1em}\setlength{\parskip}{2bp}
\renewcommand{\normalsize}{\fontsize{9.05bp}{12.1bp}\selectfont}
\ctexset{section={format=\sffamily\bfseries\fontsize{11bp}{14bp}\selectfont,beforeskip=10bp,afterskip=4bp},
  subsection={format=\sffamily\bfseries\fontsize{9.6bp}{12.1bp}\selectfont,beforeskip=8bp,afterskip=3bp}}
\fancyhead[L]{\sffamily\bfseries\fontsize{9bp}{11bp}\selectfont 观点文章／分析文章（示例）}
\fancyhead[R]{\sffamily\fontsize{8bp}{10bp}\selectfont DOI：按原稿核验}
\renewcommand{\headrulewidth}{.3bp}\renewcommand{\footrulewidth}{.3bp}
\fancyfoot[L]{\sffamily\fontsize{7bp}{9bp}\selectfont Nature Methods｜卷期与原刊页码待核验}
\fancyfoot[R]{\sffamily\fontsize{7bp}{9bp}\selectfont 非官方中性样板｜中文译页\ \thepage}
\newcommand{\MethodsCite}[1]{\textsuperscript{\hyperlink{ref:#1}{#1}}}
\newenvironment{methodscols}{\begin{multicols}{2}}{\end{multicols}}
\newcommand{\NeutralBody}{%
本段为原创版式演示文字，不陈述研究结论。正式译稿应逐段对应原文，保留比较对象、限定条件、数字、引文与否定关系。若原文跨栏或跨页续句，应先恢复阅读顺序，再组织中文段落。

版面与内容应分别维护。中文长度可以改变分页，但不应改变内容边界。方法、图注、表格脚注和声明各有独立作用，不能因文字相似而合并，也不能把未取得的附件算入已完成范围。

此处仅示范蓝色数字上标及书目跳转\MethodsCite{1}。参考条目是明确标识的中性占位，使用样板时必须以目标原稿中的完整条目替换，不沿用示例作者、出版日期或条目数量。
}
\begin{document}\normalsize
\fancypagestyle{methodsfirst}{\fancyhead{}\renewcommand{\headrulewidth}{0pt}}
\thispagestyle{methodsfirst}
{\color{MethodsOrange}\hrule height 3bp}\vspace{3bp}
{\bfseries\fontsize{23bp}{27bp}\selectfont nature methods\par}\vspace{23bp}
\noindent\sffamily\textbf{观点文章／分析文章（示例）}\hfill
{\fontsize{8bp}{10bp}\selectfont DOI：按原稿核验}\par
\vspace{4bp}\hrule height .3bp\vspace{7bp}
{\normalfont\bfseries\fontsize{25bp}{31bp}\selectfont 中文题名占位：从内容边界到完整呈现\par}\vspace{17bp}
\noindent\begin{minipage}[t]{.318\textwidth}\vspace{0pt}\sffamily\fontsize{8bp}{11bp}\selectfont
\hrule\vspace{5bp}收稿：按原稿填写\par\vspace{5bp}\hrule\vspace{5bp}
接受：按原稿填写\par\vspace{5bp}\hrule\vspace{5bp}
在线发表：按原稿填写\par\vspace{5bp}\hrule\vspace{8bp}
非官方中文译稿样板。\par 仅供学术交流。\par 出版物版权归原作者及出版商所有。
\end{minipage}\hfill\begin{minipage}[t]{.658\textwidth}\vspace{0pt}
{\sffamily\bfseries\fontsize{9bp}{12bp}\selectfont 示例作者甲\textsuperscript{1}、示例作者乙\textsuperscript{2}\par}
\vspace{8bp}\hrule\vspace{7bp}\normalfont\fontsize{10.2bp}{14bp}\selectfont
本段为中性摘要占位，演示题名、日期窄栏与作者摘要宽栏的层级关系，不陈述真实研究发现。正式译稿应完整呈现原文的问题、方法、结果及限定条件，而不是用内容总结替代全文翻译。原稿未填写的在线日期应如实标注，不能从文件名、接受日期或标识符推断。摘要下方接入双栏正文；全文中的表格、框文和图注，应按实际尺寸安排。中文长度会改变分页，本样板不承诺与英文逐页相同。
\end{minipage}\par\vspace{12bp}
\noindent{\sffamily\fontsize{7.2bp}{9.6bp}\selectfont
\textsuperscript{1}示例作者单位一（中性占位）。
\textsuperscript{2}示例作者单位二（中性占位）。通信信息按原稿填写。}\par
\vspace{8bp}\hrule\vspace{8bp}\normalfont\normalsize
\begin{methodscols}
\section*{正文结构示例}
\NeutralBody
\subsection*{版本与出处}
已出版版本和未填写在线日期的稿件，可使用同一结构，但出版字段必须依各自来源核对。书目信息缺失时明确缺项，不能以相似论文的信息补齐。此处文字仅用于展示层级与正常中文断行。
\end{methodscols}

\clearpage
\section*{通栏表格与框文}
\noindent{\sffamily\bfseries 表1｜内容对象与校核动作（原创中性示例）}\par\vspace{4bp}
{\sffamily\fontsize{8.3bp}{11.3bp}\selectfont\setlength{\tabcolsep}{7bp}
\noindent\begin{tabular}{@{}p{.23\textwidth}p{.35\textwidth}p{.35\textwidth}@{}}
\toprule
对象 & 需要保留的内容 & 校核动作\\\midrule
摘要 & 原文问题、结果与限定 & 与对应原段逐句核对\\[5bp]
正文 & 连续论述及引文关系 & 检查跨栏、跨页续句\\[5bp]
表格 & 每一行、分组与脚注 & 核对列名、单位和说明\\[5bp]
图像 & 完整面板及图注 & 逐幅检查边界与清晰度\\[5bp]
声明 & 作者、资助与许可信息 & 检查内容边界与缺项\\[5bp]
出处 & 版本所载出版字段 & 区分未知字段与已核实字段\\
\bottomrule
\end{tabular}}
\par\vspace{4bp}{\sffamily\fontsize{7.7bp}{10.3bp}\selectfont
本表只演示通栏中文表及表脚注，不是特定期刊的固定报告要求。行数和字段均为原创占位；正式长表应按需要使用跨页表格并重复列头。}\par\vspace{15bp}
\begingroup\setlength{\fboxsep}{10bp}
\noindent\colorbox{MethodsBox}{\begin{minipage}{\dimexpr\textwidth-20bp\relax}
\sffamily{\bfseries 框1｜独立内容单元（原创中性示例）}\par\smallskip
\normalfont\fontsize{8.6bp}{11.6bp}\selectfont
框文可能承担定义、讨论或方法说明，其位置与正文不同，但仍属于全文覆盖范围。应先确认框内阅读顺序，并把它与框外的续段区分开。框文在双栏页面上方跨越两栏时，可以独立排成全宽模块；不要将正文中断于框边的句子误拼入框内。

这里演示浅灰底色、紧凑标题与可读正文。正式文稿应根据框文长度重新计算高度；过长内容可以另页呈现，不能为维持原页数而缩小到难以阅读的字号。
\end{minipage}}
\endgroup\par\vspace{13bp}
\begin{methodscols}
\subsection*{双栏正文的衔接}
\NeutralBody
\subsection*{版面校核}
宽表和框文之后可以继续双栏正文。实际任务应按原文安排位置，不将样板人工分页当成固定规则。图、表与对应说明应保持可追溯关系。检查提取文字不能代替逐页渲染，因为重叠、截断或图像模糊可能无法从纯文本发现。
\end{methodscols}

\clearpage
\begin{center}\begin{minipage}{.78\textwidth}\DemoVector\end{minipage}\end{center}
\noindent{\sffamily\fontsize{7.7bp}{10.3bp}\selectfont
\textbf{图1｜内容处理流程（原创中性示例）。}
\textbf{A}，输入节点表示原文的题名、摘要、正文和图表。
\textbf{B}，处理节点表示保留限定条件，并检查单位、阅读顺序及内容边界。
\textbf{C}，输出节点表示具备可追溯引文与完整出处的中文译稿。
箭头仅表示本示例的处理顺序，不代表因果效应或数量变化。图由TikZ原生绘制，无真实数据、原图或期刊标识；图和完整图注均采用通栏。}\par\vspace{11bp}
\begin{methodscols}
\section*{讨论与交付示例}
完整译稿需要同时核对内容和呈现。图内英文若保留，应附中文标签释义，并如实说明处理范围。原始输入应保留不变；源文件包只收入复现最终版面所需的干净素材及可编辑内容\MethodsCite{2}。

真实文稿中的正文图应在参考文献之前排完。下列两条文献仅示范数字上标的内部跳转，不对应真实论文。原文信息置于全部正文、书目和声明之后，作为可追溯的结束部分。
\section*{参考文献（中性占位）}
{\fontsize{7.6bp}{10bp}\selectfont\setlength{\parindent}{0pt}\interlinepenalty=10000
\hypertarget{ref:1}{}1. Example author. Example source title. Publication details to be replaced from the supplied source. （示例占位，非真实文献。）\par\smallskip
\hypertarget{ref:2}{}2. Example author. Another example source title. Publication details to be verified. （示例占位，非真实文献。）\par}
\section*{声明（中性占位）}
致谢、资助、作者贡献、利益冲突、数据及许可等声明，应逐项按原稿翻译。本文无真实研究、实验或资助；此处仅示范声明区的位置。
\section*{原文信息}
作者、英文题名、期刊、年份、卷期、页码或文章号、DOI及明确载明的在线日期，均须从目标原稿核实并填写。以上字段在本样板中均为占位，缺项不得编造。

本文件为Nature Methods风格的非官方中性样板。刊名以普通文字显示，版面依据已核验版本作中文适配，不代表官方授权，也不承诺逐页相同。仅供学术交流；出版物版权归原作者及出版商所有。
\end{methodscols}
\end{document}
